\subsection{Locally compact fields}

DEFINITION 6. --- Let K be a locally compact field \(^{(1)}\) . For \(a \in K^{*}\) , one calls modulus of a, and denotes by \(\operatorname{mod}_{\mathbf{K}}(a)\) or simply \(\operatorname{mod}(a)\) , the modulus of the automorphism \(x \mapsto ax\) of the additive group \(K^{+}\) underlying K; one sets \(\operatorname{mod}(0) = 0\) .

Examples. --- 1) Let \(\mathbf{K} = \mathbf{R}\) . If \(s > 0\) then \(s \cdot [0,1] = [0,s]\) ; if \(s < 0\) , \(s \cdot [0,1] = [s,0]\) . Thus \(\operatorname{mod}_{\mathbf{R}} t = |t|\) for all \(t \in \mathbf{R}\) .

\begin{enumerate}
\def\labelenumi{\arabic{enumi})}
\setcounter{enumi}{1}
\tightlist
\item
  Let \(\mathbf{K} = \mathbf{Q}_p\) . If \(s \in \mathbf{Q}_p^*\) is such that \(|s|_p = p^{-n}\) , then \(s\mathbf{Z}_p\) is the set of \(x \in \mathbf{Q}_p\) such that \(|x|_p \leqslant p^{-n}\) ; therefore, if \(\mu\) denotes the normalized Haar measure on \(\mathbf{Q}_p\) , then \(\mu(s\mathbf{Z}_p) = p^{-n}\) . Thus \(\operatorname{mod}_{\mathbf{Q}_p} t = |t|_p\) for all \(t \in \mathbf{Q}_p\) .
\end{enumerate}

PROPOSITION 12. --- The function mod is continuous on K, and \(\text{mod}(ab) = \text{mod}(a)\text{mod}(b)\) for all a, b in K.

The last assertion is obvious. Prop. 4 of No.~4 shows that the function mod is continuous at every point of \(K^{*}\) . It remains only to show its continuity at 0. This is obvious for discrete K; we shall therefore assume K nondiscrete. Let \(\alpha\) be a Haar measure on \(K^{+}\) and let C be a compact subset of K such that \(\alpha(\mathrm{C}) > 0\) ; for \(a \in K^{*}\) , we have \(\alpha(a\mathrm{C}) = \mathrm{mod}(a)\alpha(\mathrm{C})\) .

Since K is not discrete, \(\alpha(\{0\}) = 0\) (No.~2, Prop. 2); therefore, for every \(\varepsilon > 0\) there exists an open neighborhood U of 0 such that \(\alpha(U) \leqslant \varepsilon\) . Since the product in K is continuous, \(aC \subset U\) for a sufficiently near 0, and then \(\text{mod}(a) \leqslant \varepsilon / \alpha(C)\) .

PROPOSITION 13. --- For every M \textgreater{} 0, let \(V_{M}\) be the set of \(x \in K\) such that \(\text{mod}(x) \leqslant M\) . If K is nondiscrete, the \(V_{M}\) form a fundamental system of compact neighborhoods of 0 in K.

The \(V_{M}\) are closed neighborhoods of 0 by Prop. 12. Let us show that they are compact. Let U be a compact neighborhood of 0. There exists an \(r \neq 0\) in K such that \(\text{mod}(r) < 1\) and \(r^{n} \in U\) for all n \textgreater{} 0: for, let W be a neighborhood of 0 such that \(WU \subset U\) ; by Prop. 12, there exists an \(r \neq 0\) in K such that \(\text{mod}(r) < 1\) and \(r \in U \cap W\) ; then \(r^{2} \in WU \subset U\) , and \(r^{n} \in U\) for all n \textgreater{} 0 by induction on n.~We are going to show that \(V_{M}\) is contained in a finite union of sets \(r^{-q}U\) (q an integer \(\geqslant 0\) ), which will prove that the \(V_{M}\) are indeed compact. If x is a cluster point of the sequence ( \(r^{n}\) ), then \(\text{mod}(x)\) is a cluster point of the sequence ( \(\text{mod}(r)^{n}\) ), therefore \(\text{mod}(x) = 0\) , x = 0; since U is compact, it follows (GT, I, §9, No.~1, Cor. of Th. 1) that \(\lim_{n \to \infty} r^{n} = 0\) . Now let \(a \in V_{M}\) . Since the sequence ( \(r^{n}a\) ) \(_{n \geqslant 0}\) tends to 0, there exists a smallest integer \(n \geqslant 0\) such that \(r^{n}a \in U\) . If n \textgreater{} 0 then \(r^{n-1}a \notin U\) , therefore \(r^{n}a \in U \cap C(rU)\) ; the closure X of \(U \cap C(rU)\) is compact since U is compact, and it does not contain 0 since rU is a neighborhood of 0; therefore, in X, \(\text{mod}(x)\) is bounded below by a number m \textgreater{} 0. Thus, if n \textgreater{} 0, we have \(m \leqslant \text{mod}(r^{n}a)\) , whence \(\text{mod}(r^{-1})^{n} \leqslant M/m\) . Since \(\text{mod}(r^{-1}) > 1\) , the integer n can only take on a finite number of values, a number not depending on a, which completes the proof of our assertion.

This being so, since the intersection of the \(\mathrm{V_M}\) reduces to \(\{0\}\) the \(\mathrm{V_M}\) form a fundamental system of neighborhoods of 0 (GT, I, §9, No.~2, Prop. 1).

COROLLARY. --- The topology of a nondiscrete locally compact field admits a countable base.

For, K is the union of compact sets \(V_{1}, V_{2}, \ldots\) . On the other hand, K is metrizable by Prop. 1 of GT, IX, §3, No.~1. Therefore the topology of K admits a countable base (loc. cit., §2, No.~9, Cor. of Prop. 16).

PROPOSITION 14. --- Let \(\alpha\) be a Haar measure on \(\mathbf{K}^{+}\) . Then the measure \(\beta = (\mathrm{mod}_{\mathbf{K}})^{-1} \cdot \alpha\) on \(\mathbf{K}^{*}\) is a left Haar measure on the multiplicative group \(\mathbf{K}^{*}\) .

For, if \(b \in \mathbf{K}^*\) , the mapping \(a \mapsto b^{-1}a\) of \(\mathbf{K}\) into \(\mathbf{K}\) transforms \(\alpha\) into \((\mathrm{mod}_{\mathbf{K}}b)\alpha\) , hence \((\mathrm{mod}_{\mathbf{K}})^{-1} \cdot \alpha\) into itself, whence the proposition.

COROLLARY. --- Let f be a function defined on \(K^{*}\) , with values in \(\overline{R}\) or in a Banach space. For f to be \(\beta\) -integrable, it is necessary and sufficient that \((\mathrm{mod}_K)^{-1}f\) be \(\alpha\) -integrable, in which case

\[
\int_ {\mathrm{K} ^ {*}} f (x) d \beta (x) = \int_ {\mathrm{K} ^ {+}} \left(\mathrm{mod} _ {\mathrm{K}} (x)\right) ^ {- 1} f (x) d \alpha (x).
\]

This follows from Prop. 14, the Cor. of Prop. 13, and Ch. V, §5, No.~3, Th. 1.

PROPOSITION 15. --- Assume K to be commutative. Let u be an automorphism of the vector space \(E = K^{n}\) . Then

\[
\operatorname{mod} _ {\mathbf {E}} u = \operatorname{mod} _ {\mathbf {K}} (\det u).
\]

It suffices to verify the formula when \(u\) runs over a system of generators of \(\mathbf{GL}(\mathbf{E})\) . Now, \(\mathbf{GL}(\mathbf{E})\) is generated by the following elements (A, II, §10, No.~13, Cor. 2 of Prop. 14):

\begin{enumerate}
\def\labelenumi{(\alph{enumi})}
\tightlist
\item
  The elements \(u_{1}\) of the form
\end{enumerate}

\[
\left(x _ {1}, \dots , x _ {n}\right) \mapsto \left(x _ {\sigma (1)}, \dots , x _ {\sigma (n)}\right),
\]

where \(\sigma \in \mathfrak{S}_n\) ;

\begin{enumerate}
\def\labelenumi{(\alph{enumi})}
\setcounter{enumi}{1}
\tightlist
\item
  the elements \(u_{2}\) of the form
\end{enumerate}

\[
\left(x _ {1}, \dots , x _ {n}\right) \mapsto \left(a x _ {1}, x _ {2}, \dots , x _ {n}\right)
\]

with \(a\in \mathbf{K}^*\)

\begin{enumerate}
\def\labelenumi{(\alph{enumi})}
\setcounter{enumi}{2}
\tightlist
\item
  the elements \(u_{3}\) of the form
\end{enumerate}

\[
\left(x _ {1}, \dots , x _ {n}\right) \mapsto \left(x _ {1} + \sum_ {i = 2} ^ {n} c _ {i} x _ {i}, x _ {2}, \dots , x _ {n}\right).
\]

If \(f \in \mathcal{K}(\mathbf{E})\) then, denoting by \(\alpha\) a Haar measure on \(\mathbf{K}^+\) ,

\[
\begin{array}{l} \int \dots \int_ {\mathrm{K} ^ {n}} f (x _ {1} + \sum_ {i = 2} ^ {n} c _ {i} x _ {i}, x _ {2}, \ldots , x _ {n}) d \alpha (x _ {1}) d \alpha (x _ {2}) \ldots d \alpha (x _ {n}) \\ \quad = \int \dots \int_ {\mathrm{K} ^ {n - 1}} d \alpha (x _ {2}) \ldots d \alpha (x _ {n}) \int_ {\mathrm{K}} f (x _ {1} + \sum_ {i = 2} ^ {n} c _ {i} x _ {i}, x _ {2}, \ldots , x _ {n}) d \alpha (x _ {1}) \\ \quad = \int \dots \int_ {\mathrm{K} ^ {n - 1}} d \alpha (x _ {2}) \ldots d \alpha (x _ {n}) \int_ {\mathrm{K}} f (x _ {1}, x _ {2}, \ldots , x _ {n}) d \alpha (x _ {1}) \\ \quad = \int \dots \int_ {\mathrm{K} ^ {n}} f (x _ {1}, \ldots , x _ {n}) d \alpha (x _ {1}) \ldots d \alpha (x _ {n}), \end{array}
\]

and on the other hand \(\mathrm{mod}_{\mathrm{K}}(\det u_{3})=\mathrm{mod}_{\mathrm{K}}(1)=1\) , whence the result for \(u_{3}\) . It is established in an analogous manner for \(u_{1}\) and \(u_{2}\) .

Let K be a commutative locally compact field, E a vector space of finite dimension n over K. If \(\varphi\) is an isomorphism of the vector space \(K^{n}\) onto the vector space E, \(\varphi\) transforms the topology of \(K^{n}\) into a topology on E that makes E a locally compact vector space. This topology (called canonical) is independent of \(\varphi\) since every automorphism of the vector space \(K^{n}\) is bicontinuous. Absent express mention to the contrary, when we speak of E as a topological vector space, the topology will always be understood to be the one just defined. Every automorphism u of the vector space E is bicontinuous, therefore \(mod_{E}u\) is defined. If, on the other hand, u is a noninvertible endomorphism of E, one sets \(mod_{E}u = 0\) . Then:

COROLLARY 1. --- Let K be a commutative locally compact field, E a finite-dimensional vector space over K, and u an endomorphism of the vector space E. Then \(\text{mod}_{\mathbf{E}}(u) = \text{mod}_{\mathbf{K}}(\det u)\) .

If \(u\) is invertible, this follows from Prop. 15. If \(u\) is not invertible then \(\det u = 0\) , therefore \(\operatorname{mod}_{\mathbf{K}}(\det u) = 0 = \operatorname{mod}_{\mathbf{E}}u\) .

COROLLARY 2. --- Let E be a real vector space of finite dimension n, \((e_{1}, e_{2}, \ldots, e_{n})\) a basis of E, P the set of \(x = \sum_{i=1}^{n} \xi_{i} e_{i} \in E\) such that \(0 \leqslant \xi_{i} \leqslant 1\) for all i, and \(\mu\) the unique Haar measure on the additive group E such that \(\mu(P) = 1\) . Let \(x_{1}, \ldots, x_{n}\) be points of E, S the closed convex envelope in E of the set \(\{0, x_{1}, \ldots, x_{n}\}\) . Writing \(x_{i} = \sum_{j=1}^{n} \alpha_{ij} e_{j}\) , one has

\[
\mu (\mathrm{S}) = \mu (\stackrel {\circ} {\mathrm{S}}) = \frac {1}{n !} | \det (\alpha_ {i j}) |.
\]

We shall identify E with \(R^{n}\) by means of the isomorphism that transforms \((e_{i})\) into the canonical basis of \(R^{n}\) . Then \(\mu\) is identified with the Lebesgue measure \(\mu_{n}\) on \(R^{n}\) .

Suppose first that \(x_{i} = e_{i}\) for all \(i\) . Then \(S\) is the set \(S_{n}\) of \(x = (\xi_{i})\) in \(\mathbf{R}^n\) such that

\[
\xi_ {i} \geqslant 0 \text {   for   all   } i \quad \text { and } \quad \xi_ {1} + \dots + \xi_ {n} \leqslant 1.
\]

Set \(\mu_n(S_n) = a_n\) . Let \(\lambda \in \mathbf{R}\) . Identifying \(\mathbf{R}^n\) with \(\mathbf{R}^{n-1} \times \mathbf{R}\) , we may consider the section \(C_\lambda\) of \(S_n\) at \(\lambda\) . This section is empty if \(\lambda < 0\) or \(\lambda > 1\) ; if \(0 \leqslant \lambda \leqslant 1\) , \(C_\lambda\) is the set of \((\xi_1, \ldots, \xi_{n-1}) \in \mathbf{R}^{n-1}\) such that

\[
\xi_ {1} \geqslant 0, \dots , \xi_ {n - 1} \geqslant 0, \quad \xi_ {1} + \dots + \xi_ {n - 1} \leqslant 1 - \lambda ,
\]

hence may be deduced from \(S_{n-1}\) by a homothety of ratio \(1-\lambda\) , so that \(\mu_{n-1}(C_{\lambda})=(1-\lambda)^{n-1}a_{n-1}\) . By the Lebesgue--Fubini theorem,

\[
a _ {n} = \int_ {0} ^ {1} (1 - \lambda) ^ {n - 1} a _ {n - 1} d \lambda = \frac {1}{n} a _ {n - 1}.
\]

Since \(a_1 = 1\) , one sees that \(a_n = \frac{1}{n!}\) .

Let us return to the general case of the corollary. Let u be the endomorphism of \(R^{n}\) such that \(u(e_{i}) = x_{i}\) for all i. One has \(u(S_{n}) = S\) . If u is invertible, Prop. 15 proves that

\[
\mu_ {n} (\mathrm{S}) = \frac {1}{n !} | \det u | = \frac {1}{n !} | \det (\alpha_ {i j}) |.
\]

Since \(S - \mathring{S}\) is contained in a finite number of hyperplanes, \(\mu(\mathring{S}) = \mu(S)\) . Finally, if \(u\) is not invertible, then \(S\) is contained in a hyperplane, so that \(\mu(S) = 0 = \det(\alpha_{ij})\) .

